\oldpage[II]{50}
graduée (\hyperref[I.1.4.3-fr]{\textbf{I}, 1.4.3})~; posons
$X_U=\operatorname{Proj}(\Gamma(U,\mathcal{S}))$. Soient $U'\subset U$ un
second ouvert affine de $Y$, $A'$ son anneau, $j:U'\to U$ l'injection
canonique, qui correspond à l'homomorphisme de restriction $A\to A'$~; on a
$\mathcal{S}|U'=j^*(\mathcal{S}|U)$, et par suite
$S'=\Gamma(U',\mathcal{S})$ s'identifie canoniquement à $S\otimes_A A'$
(\hyperref[I.1.6.4-fr]{\textbf{I}, 1.6.4}). On en conclut
(\hyperref[II.2.8.10-fr]{2.8.10}) que $X_{U'}$ s'identifie canoniquement à
$X_U\times_U U'$, et par suite aussi à $f_U^{-1}(U')$, en désignant par
$f_U$ le morphisme structural $X_U\to U$
(\hyperref[I.4.4.1-fr]{\textbf{I}, 4.4.1}). Nous désignerons par $\sigma_{U',U}$
l'isomorphisme canonique $f_U^{-1}(U')\xrightarrow{\sim}X_{U'}$ ainsi défini,
par $\rho_{U',U}$ l'immersion ouverte $X_{U'}\to X_U$ obtenue en composant
$\sigma_{U',U}^{-1}$ et l'injection canonique $f_U^{-1}(U')\to X_U$. Il est
immédiat que si $U''\subset U'$ est un troisième ouvert affine de $Y$, on a
$\rho_{U'',U}=\rho_{U'',U'}\circ\rho_{U',U}$.
\end{env}

\begin{proposition}[3.1.2]
\label{II.3.1.2-fr}
Soit $Y$ un préschéma. Pour toute $\mathcal{O}_Y$-Algèbre graduée
quasi-cohérente $\mathcal{S}$ à degrés positifs, il existe un préschéma $X$
au-dessus de $Y$, et un seul à un $Y$-isomorphisme près, ayant la propriété
suivante~: si $f$ est le morphisme structural $X\to Y$, pour tout ouvert
affine $U$ de $Y$, il existe un isomorphisme $\eta_U$ du préschéma
induit $f^{-1}(U)$ sur $X_U=\operatorname{Proj}(\Gamma(U,\mathcal{S}))$ tel
que, si $V$ est un second ouvert affine de $Y$ contenu dans $U$, le diagramme
\[
\label{II.3.1.2.1-fr}
  \xymatrix{
    f^{-1}(V)\ar[r]^{\overset{\eta_V}{\sim}}\ar[d]
    & X_V\ar[d]^{\rho_{V,U}}\\
    f^{-1}(U)\ar[r]^{\sim}_{\eta_U}
    & X_U
  }
\tag{3.1.2.1}
\]
soit commutatif.
\end{proposition}

\begin{proof}
Pour deux ouverts affines $U$, $V$ de $Y$, soit $X_{U,V}$ le préschéma
induit sur $f_U^{-1}(U\cap V)$ par $X_U$~; nous allons définir un
$Y$-isomorphisme $\theta_{U,V}:X_{V,U}\xrightarrow{\sim}X_{U,V}$. Pour cela,
considérons un ouvert affine $W\subset U\cap V$~: en composant les
isomorphismes
\[
  f_U^{-1}(W)\xrightarrow{\sigma_{W,U}}X_W
  \xrightarrow{\sigma_{W,V}^{-1}}f_V^{-1}(W),
\]
on obtient un isomorphisme $\tau_W$, et on vérifie aussitôt que si
$W'\subset W$ est un ouvert affine, $\tau_{W'}$ est la restriction à
$f_U^{-1}(W')$ de $\tau_W$~; les $\tau_W$ sont donc bien les restrictions
d'un $Y$-isomorphisme $\theta_{V,U}$. En outre, si $U$, $V$, $W$ sont trois
ouverts affines de $Y$, $\theta'_{U,V}$, $\theta'_{V,W}$ et
$\theta'_{U,W}$ les restrictions de $\theta_{U,V}$, $\theta_{V,W}$,
$\theta_{U,W}$ aux images réciproques de $U\cap V\cap W$ dans $X_V$, $X_W$,
$X_W$ respectivement, il résulte des définitions précédentes que l'on a
$\theta'_{U,V}\circ\theta'_{V,W}=\theta'_{U,W}$. L'existence de $X$
vérifiant les propriétés énoncées résulte donc de
(\hyperref[I.2.3.1-fr]{\textbf{I}, 2.3.1})~; son unicité à un $Y$-isomorphisme près
est triviale, compte tenu de (\hyperref[II.3.1.2.1-fr]{3.1.2.1}).
\end{proof}

\begin{env}[3.1.3]
\label{II.3.1.3-fr}
Nous dirons que le préschéma $X$ défini dans
(\hyperref[II.3.1.2-fr]{3.1.2}) est le \emph{spectre homogène} de la
$\mathcal{O}_Y$-Algèbre graduée quasi-cohérente $\mathcal{S}$ et nous le
noterons $\operatorname{Proj}(\mathcal{S})$. Il est immédiat que
$\operatorname{Proj}(\mathcal{S})$ est \emph{séparé au-dessus de $Y$}
((\hyperref[II.2.4.2-fr]{2.4.2}) et
(\hyperref[I.5.5.5-fr]{\textbf{I}, 5.5.5}))~; si $\mathcal{S}$ est une
$\mathcal{O}_Y$-Algèbre \emph{de type fini}
(\hyperref[I.9.6.2-fr]{\textbf{I}, 9.6.2}) $\operatorname{Proj}(\mathcal{S})$ est
\emph{de type fini} sur $Y$ ((\hyperref[II.2.7.1-fr]{2.7.1, (ii)}) et
(\hyperref[I.6.3.1-fr]{\textbf{I}, 6.3.1})).

Si $f$ est le morphisme structural $X\to Y$, il est immédiat que pour tout
préschéma induit par $Y$ sur un ouvert $U$ de $Y$, $f^{-1}(U)$ s'identifie au
spectre homogène $\operatorname{Proj}(\mathcal{S}|U)$.
\end{env}

\begin{proposition}[3.1.4]
\label{II.3.1.4-fr}
Soit $f\in\Gamma(Y,\mathcal{S}_d)$ pour $d>0$. Il existe une partie ouverte
$X_f$ de l'espace sous-jacent à $X=\operatorname{Proj}(\mathcal{S})$ ayant la
propriété suivante~: pour tout ouvert affine $U$ de $Y$,
